\section{General Discussion}

Across two studies, we probed the prosocial premise of criminal punishment. Study~1 revealed that punitiveness is embedded in a dense, cohesive psychological mindset in which hostile aggression ($r = .59$) dramatically outpredicts crime concerns ($r = .32$). Study~2 showed that despite this hostile substrate, participants overwhelmingly invoked prosocial language when justifying their sentences---deterrence, public safety, and rehabilitation. The prosocial surface, however, did not track the psychological substrate. Of the 33 cells in the facade matrix that reached both significance and a meaningful effect size, 32 were negative. High-hostile participants did not endorse more revenge or suffering. They quietly dropped one specific prosocial theme, rehabilitation, while continuing to invoke deterrence, incapacitation, and proportional justice at rates indistinguishable from low-hostile participants. The prosocial framing functions as a cultural default with selective suppression: a shared rhetorical script available to liberals and conservatives, hostile and non-hostile, with the rehabilitation theme attenuated specifically among those highest in punitive dispositions.

\subsection{The Punitiveness Mindset}

Recall that concern about crime provides the logical and normatively appropriate impetus for criminal punishment \citep{Enns2016, Pickett2019, Roberts2003}. As such, the correlation between the crime concerns cluster and punitiveness is expected to be positive and strong. We found the correlation was indeed positive ($r = .32$), but it was the weakest of the four clusters. We find it striking that the most prosocially justifiable correlate is also the weakest predictor.

The emotions cluster correlated fairly strongly with punitiveness ($r = .53$). Among the constituent emotions, punitiveness correlated most strongly with hatred of criminals ($r = .53$). Unlike anger, hatred is a stable feeling that precedes and transcends particular events. Hatred has been labeled ``our most dangerous emotion'' \citep{Brogaard2020} and is widely recognized as triggering the desire to coerce, exclude, harm, and humiliate its object \citep{BarTal2007, FischerHalperinCanettiJasini2018}.

Punitiveness was also positively associated with measures of personality and ideology ($r = .47$). The most prominent within this cluster were right-wing authoritarianism ($r = .50$) and racial resentment ($r = .46$), with social dominance orientation also significant ($r = .32$). Punishment decisions thus appear to reflect the punisher's ideological and attitudinal commitments---particularly troubling with respect to the anti-egalitarian sentiments embodied by SDO and the racial prejudice embodied by racial resentment. These findings extend a lineage of research highlighting the intimate connection between racism and criminal punishment \citep{Armour2020, BoboJohnson2004, ChiricosWelchGertz2004, UnneverCullen2010a}.

Perhaps the starkest finding is the prominence of the hostile aggression cluster, with a correlation of $r = .59$. Punitiveness was closely interrelated with hating people convicted of crimes, degrading them, inflicting suffering on them, tolerating physical and sexual assaults on them in prison, and exacting revenge on them. These correlates resonate with Garfinkel's \citeyearpar{Garfinkel1956} view of punishment as status degradation, and they complicate Durkheim's justification of punishment as a promoter of social solidarity \citep[see][]{Garland1990}: solidarity for some is purchased at the exclusion of others. George Herbert Mead \citeyearpar{Mead1918} proposed that such solidarity is best conceived as a ``solidarity of aggression''---a conception recently advanced as ``hostile solidarity'' by \citet{Carvalho2018} and \citet{ChamberlenCarvalho2019}.

\subsection{Goal Mischaracterization}

How do we reconcile the prosociality that dominates the punishment discourse with the hostile aggression expressed by our participants? Study~2 suggests that the avowed prosociality amounts to a mischaracterization of people's true motives. Several mechanisms plausibly contribute.

First, societal norms have shifted over the centuries to soften the overtness of state violence \citep{Elias1969, Foucault1977, Garland1991}. People in general harbor a deep-seated aversion to acting aggressively \citep{CrockettEtAl2014, CushmanEtAl2012} and would tend to bristle at the suggestion that they are motivated by hostile aggression. Indeed, when being observed by others, people make a special effort to behave in non-cruel ways \citep{FincherTetlock2015, GantmanEtAl2017}.

Second, people are typically unaware of the psychological forces driving their punishment preferences. Moral judgments derive from quick and inaccessible intuitions rather than from deliberate reasoning \citep{Haidt2001}. When pressed to justify those intuitions, people often fall back on bare restatements---moral dumbfounding \citep[][p.~817]{Haidt2001}. Similar findings have been reported in the punishment context. People habitually endorse deterrence and incapacitation, yet they continue to punish just as harshly when those justifications cannot plausibly apply, and they tend to stick to their positions even when pressed on the implausibility \citep{AharoniFridlund2012, Carlsmith2008, CarlsmithDarley2008, CarlsmithEtAl2002}.

Our NLP findings provide novel and direct empirical evidence for this account. The computational text analysis quantifies the discrepancy between the prosocial language of punishment and the hostile psychology that drives the underlying attitudes. Crucially, our multiverse analysis indicates that this discrepancy is not borne at the individual level. Four of the six prosocial--dark gap measures returned cultural-default verdicts: no relationship between language and hostile aggression. The two measures that did show individual-level effects (Claude $r = -.25$, DeBERTa $r = -.17$) were modest and concentrated in the suppression of the rehabilitation theme rather than the open expression of dark themes. The language of deterrence and public safety is not deployed strategically by hostile people to mask their motives. Rather, it is the shared vocabulary that everyone uses when talking about punishment, including the people whose actual psychological profile bears little relation to that vocabulary.

What differs across the hostility distribution is not whether prosocial language is invoked but which prosocial theme is invoked. The most punitive and hostile participants quietly dampen rehabilitation-themed language while continuing to use deterrence, incapacitation, and proportional justice at the same rate as everyone else. They do not openly endorse revenge, suffering, or exclusion in their language. The hostile substrate that Study~1 documented is invisible on the linguistic surface---except through the silence around rehabilitation.

\subsection{Implications}

Several implications follow from these findings.

First, they challenge the taken-for-granted prosocial framing of punishment. The fact that hostile aggression---not concern about crime---is the dominant correlate of punitiveness suggests that much of what fuels public demand for harsh sentences has little to do with making society safer and much to do with satisfying darker impulses. The fact that these impulses are dressed in prosocial language makes them harder to identify, scrutinize, and resist.

Second, the cultural-default finding suggests that the problem is not one of individual deception. People are not consciously lying when they cite deterrence, and they may even experience those beliefs as genuine. Rather, they are drawing on a culturally available explanatory framework that has become so pervasive that it functions as the natural way to talk about punishment. This has parallels to what \citet{Wilson2004} described as the ``adaptive unconscious''---the gap between the reasons we give for our behavior and the actual causes of that behavior \citep[see also][]{Haidt2001, NisbettWilson1977}.

Third, computational text analysis captures something about punitiveness that conventional self-report scales miss. In our hierarchical regressions, text features added $\Delta R^2 = .079$ above the Study~1 psychological predictors when predicting attitudinal punitiveness, and $\Delta R^2 = .134$ when predicting the sentencing decision itself. This points to the potential of computational text analysis as a complement to structured measurement in the study of legal attitudes.

Fourth, these findings are relevant to the ongoing debate about punishment policy, especially regarding the American phenomenon known as mass incarceration \citep[see][]{Garland2020, NRC2014, Tonry2018}. If punitive policy is driven in substantial part by hostile aggression rather than legitimate crime concerns, then the usual policy debates---which take the espoused prosocial goals at face value and argue about whether punishment meets its avowed utilitarian goals---are missing a significant part of the picture. A more honest discourse would acknowledge the role of retaliatory and aggressive motivations and critically examine their contribution to punishment policy.

\subsection{Limitations}

Several limitations bear on the interpretation of these findings.

First, the study is cross-sectional and correlational. A critic could protest that correlational relationships say little about the causal nature of the interaction. Such backward reasoning is theoretically possible \citep{HolyoakSimon1999, SimonRead2023}, but with the exception of anger, most correlates capture stable features of the person---hatred, trait aggression, SDO, racial resentment---that are unlikely to be swayed substantially by a particular engagement with punishment \citep{JuddKrosnick1982, Skitka2010}.

Second, the theoretical clustering was guided by face validity rather than confirmed by the data's factor structure. CFA yielded fit indices below conventional thresholds ($\chi^2(98) = 699.10$, CFI $= .86$, RMSEA $= .11$), indicating that the clusters are best understood as conceptual organizers rather than as a strict latent measurement model. The substantive findings, which operate at the level of observed cluster composites, are not affected by the measurement-model limitation. Sensitivity analyses further demonstrated that the hostile aggression advantage survived even the strictest specifications, with all three stringency-level analyses ranging from $\Delta r = .22$ to $\Delta r = .26$ and all 95\% bootstrap CIs excluding zero.

Third, our sample was recruited through the Prolific pool of volunteer participants. Although we used a nationally representative quota sample, online samples may differ from the broader population in ways not captured by quota variables \citep{ChandlerShapiro2016}. Social desirability \citep{Fisher1993} likely led participants to underestimate their true levels on the harsh measures, raising the possibility that the true motivations animating punitiveness are even darker than we observed.

Fourth, the NLP methods depend on specific prototype sentences and classification labels. Sensitivity analyses showed robustness across five prototype sets and three embedding models (45 combinations in total), but specific point estimates should be interpreted with appropriate caution.

Fifth, the open-ended prompt may have encouraged post-hoc rationalization. Future studies could elicit justifications before the sentencing decision, or use think-aloud protocols, to examine whether the cultural default operates even when participants are not constructing retrospective explanations.

Sixth, the sentence recommendations that we incorporated into the punitiveness measures all involved second-degree murder convictions. We increased generalizability by randomly alternating among three such vignettes, but generalizability to other crime types remains an open question.

\subsection{Conclusion}

Is criminal punishment prosocial? Our findings suggest that the answer is far more complicated than the official discourse implies. The psychological forces most strongly associated with punitiveness---hatred, revenge, degradation, suffering, authoritarianism---are orthogonal to the prosocial goals of deterrence, public safety, and rehabilitation. When people explain their punishment preferences, they reliably invoke those prosocial goals. But the prosocial language they use does not track the psychological measures that most strongly predict their punitive attitudes. The most punitive participants are not the ones who sound most prosocial, nor the ones who openly express dark themes. They are the ones who quietly stop talking about rehabilitation. This pattern is best described not as individual-level deception but as a cultural default with selective suppression: a shared rhetorical script that obscures the darker motivational substrate of punishment, and that the most punitive among us deploy with one specific theme attenuated.

Coming to terms with this complexity might lead us to reconsider how we think, talk, and act in the domain of criminal punishment. At a minimum, these findings call for a discourse that demystifies the prevailing prosocial rhetoric and makes room for a frank recognition that we punish, in part, to satisfy our own---often unsavory---psychic needs.